% 編譯方式： xelatex dp-image.tex（執行兩次以解析交叉引用）
\documentclass[a4paper,11pt,fontset=none]{ctexart}

% ---- 繁體中文字型：依序嘗試常見的繁體字型 -------------------------------
\usepackage{xeCJK}
\IfFontExistsTF{Noto Serif CJK TC}{%
  \setCJKmainfont{Noto Serif CJK TC}\setCJKsansfont{Noto Sans CJK TC}\setCJKmonofont{Noto Sans CJK TC}}{%
\IfFontExistsTF{Source Han Serif TC}{%
  \setCJKmainfont{Source Han Serif TC}\setCJKsansfont{Source Han Sans TC}\setCJKmonofont{Source Han Sans TC}}{%
\IfFontExistsTF{PMingLiU}{%
  \setCJKmainfont{PMingLiU}\setCJKsansfont{Microsoft JhengHei}\setCJKmonofont{MingLiU}}{%
\IfFontExistsTF{AR PL UMing TW}{%
  \setCJKmainfont{AR PL UMing TW}\setCJKsansfont{AR PL UMing TW}\setCJKmonofont{AR PL UMing TW}}{%
  \setCJKmainfont{Noto Serif CJK TC}}}}}

\usepackage[margin=2.5cm]{geometry}
\usepackage{amsmath,amssymb,amsthm}
\usepackage{booktabs}
\usepackage{enumitem}
\usepackage{tikz}
\usetikzlibrary{arrows.meta,positioning}
\usepackage{float}
\usepackage{algorithm}
\usepackage{algpseudocode}
\usepackage{listings}
\usepackage[hidelinks]{hyperref}

\newcommand{\tg}[2]{\texttt{#1\,|\,#2}}
\newcommand{\code}[1]{\texttt{\detokenize{#1}}}
\renewcommand{\emph}[1]{\textbf{#1}}
\setlength{\emergencystretch}{3em}
\newcommand{\F}{\mathbb{F}}
\newcommand{\E}{\mathcal{E}}
\newcommand{\ra}{\mathrm{root}}
\DeclareMathOperator{\var}{var}
\DeclareMathOperator{\supp}{supp}
\DeclareMathOperator{\lev}{level}
\DeclareMathOperator{\skipc}{skip}
\DeclareMathOperator{\low}{low}
\DeclareMathOperator{\high}{high}

\newtheorem{lemma}{引理}
\newtheorem{theorem}{定理}
\theoremstyle{definition}
\newtheorem{definition}{定義}

\floatname{algorithm}{演算法}
\lstset{breaklines=true,basicstyle=\ttfamily\small,columns=fullflexible,keepspaces=true,frame=single,xleftmargin=1em,xrightmargin=1em}
\algrenewcommand\algorithmicrequire{\textbf{輸入：}}
\algrenewcommand\algorithmicensure{\textbf{輸出：}}
\setlength{\parskip}{0.3em}

\title{以錯誤集合決策圖上的動態規劃判定不等價錯誤對}
\author{SPBDD 設計說明（\code{StabilizerCode::find_inequivalent_pair}，\code{PairMethod::Dp}）}
\date{}

\begin{document}
\maketitle

本文定義問題、說明方法、證明正確性、分析複雜度，並與原本的基底變換（compose）實作比較；compose 現在僅作為參考實作與基準測試之用，預設方法已改為動態規劃（DP）。第~\ref{sec:meas} 節給出量測數據，所有數字皆可由 \code{examples/pair_benchmark.cpp} 與 \code{examples/dp_vs_explicit.cpp} 重現。

%% =====================================================================
\section{問題}\label{sec:problem}

設 $S \le \mathcal P_n$ 為穩定子群（stabilizer group），有 $r$ 個獨立生成元 $g_1,\dots,g_r$，因此此碼有 $k=n-r$ 個邏輯量子位元。以不計相位的辛（symplectic）形式表示 Pauli 算符：$e=(x\mid z)\in\F_2^{2n}$，並沿用 SPBDD 的變數編號 $x_q=$ 變數 $2q$、$z_q=$ 變數 $2q+1$。辛內積記為
\[
\langle e,f\rangle = x_e\cdot z_f + z_e\cdot x_f \pmod 2 .
\]

取定邏輯算符代表元 $\bar X_1,\dots,\bar X_k,\bar Z_1,\dots,\bar Z_k\in N(S)$，滿足 $\langle\bar X_i,\bar Z_j\rangle=\delta_{ij}$ 且 $\langle\bar X_i,\bar X_j\rangle=\langle\bar Z_i,\bar Z_j\rangle=0$。定義兩個 $\F_2$-線性映射
\[
\sigma(e)=\bigl(\langle e,g_i\rangle\bigr)_{i=1}^{r}\in\F_2^{r},
\qquad
\lambda(e)=\bigl(\langle e,\bar Z_j\rangle_{j=1}^{k},\ \langle e,\bar X_j\rangle_{j=1}^{k}\bigr)\in\F_2^{2k},
\]
分別稱為\emph{症狀}（syndrome）與\emph{邏輯特徵}（logical signature），程式中對應 \code{syndrome()} 與 \code{logical_signature()}。兩者串接得到
\[
\pi=(\sigma,\lambda):\F_2^{2n}\to\F_2^{m},\qquad m=r+2k=n+k,
\]
即基底變換在「真正需要的座標」上的投影。

\begin{definition}[不等價錯誤對問題]
給定錯誤集合 $\E\subseteq\F_2^{2n}$（以 $2n$ 個變數上的 BDD 表示），判定是否存在 $e_1,e_2\in\E$ 使得 $e_1+e_2\in N(S)\setminus S$，並在存在時回傳這樣的一對。
\end{definition}

\begin{lemma}[成對判準]\label{lem:pair}
對任意 $e_1,e_2\in\F_2^{2n}$，
\[
e_1+e_2\in N(S)\setminus S\iff \sigma(e_1)=\sigma(e_2)\ \wedge\ \lambda(e_1)\neq\lambda(e_2).
\]
\end{lemma}
\begin{proof}
$N(S)=\{f:\sigma(f)=0\}$。又對 $f\in N(S)$，$f$ 可寫成 $\sum_i a_i g_i+\sum_j(c_j\bar X_j+f_j\bar Z_j)$，其 $\bar X,\bar Z$ 的係數可藉由與配對元 $\bar Z_j,\bar X_j$ 取辛內積讀出，故 $f\in S\iff\lambda(f)=0$。由 $\sigma,\lambda$ 的線性性，$\sigma(e_1+e_2)=0$ 且 $\lambda(e_1+e_2)\neq0$ 等價於上述條件。
\end{proof}

於是定義\emph{關係}
\[
G(\E)=\pi(\E)=\{(\sigma(e),\lambda(e)):e\in\E\}\subseteq\F_2^{r}\times\F_2^{2k},
\]
則存在不等價錯誤對，若且唯若某個症狀 $s$ 在 $G$ 中對應到至少兩個不同的邏輯特徵。

\begin{lemma}[$2k$ 個單一位元檢驗]\label{lem:bits}
這樣的 $s$ 存在，若且唯若存在某個特徵位元 $v\in\{1,\dots,2k\}$ 使得（視為 $s$ 的函數）
\[
\bigl(\exists\lambda.\ G\wedge\lambda_v=0\bigr)\ \wedge\ \bigl(\exists\lambda.\ G\wedge\lambda_v=1\bigr)
\]
可滿足。
\end{lemma}
\begin{proof}
兩個不同的特徵至少有一個位元不同；反之，若某症狀同時容許某一位元的兩種取值，它就有兩個不同的特徵。
\end{proof}

因此整個判定問題化為\emph{建構 $G(\E)$ 這個 $r+2k$ 個變數上的 BDD}；引理~\ref{lem:bits} 的 $2k$ 個檢驗對所有方法皆相同（程式中 \code{find_inequivalent_pair} 的 multiplicity 迴圈）。以下全部在討論這一步的代價。

%% =====================================================================
\section{直覺與完整範例：$[[5,1,3]]$ 碼、大小為 5 的錯誤集合}\label{sec:example}

這一節用一個能手算的例子，盡量不用術語，說明整個方法在做什麼。後面的第~\ref{sec:dp} 節只是把這裡的手算過程一般化。

\subsection{要回答的問題，白話版}

量子糾錯碼會對每個錯誤 $e$ 量出兩樣東西：
\begin{itemize}[leftmargin=1.5em]
\item \textbf{症狀 $\sigma(e)$}：$e$ 讓哪幾個穩定子「報警」（與它反交換）。解碼器\emph{只看得到}症狀。
\item \textbf{邏輯特徵 $\lambda(e)$}：$e$ 對被保護的那個邏輯量子位元造成什麼效果（是否翻轉邏輯 $\bar Z$ 與 $\bar X$ 的測量）。解碼器看不到它，但它決定修正對不對。
\end{itemize}
如果兩個不同的錯誤 $e_1,e_2$ 症狀相同，解碼器分不出是哪一個；若它們的邏輯特徵又不同，則不論解碼器怎麼猜，總有一種情形猜錯，而且猜錯留下的恰是一個邏輯算符 $e_1e_2$。所以我們要檢查的是：

\begin{center}
\fbox{\parbox{0.85\textwidth}{\centering 錯誤集合 $\E$ 裡，有沒有兩個錯誤\\\textbf{症狀相同、邏輯特徵不同}？}}
\end{center}

把每個錯誤的 $(\sigma,\lambda)$ 並排寫成一個「標籤」（tag）。問題就變成：標籤集合裡有沒有兩個標籤，前半（$\sigma$）一樣、後半（$\lambda$）不一樣？

\subsection{$[[5,1,3]]$ 碼與標籤表}

五量子位元碼有 $r=4$ 個穩定子、$k=1$ 個邏輯量子位元，因此標籤是 $4+2=6$ 個位元，寫成 \tg{$\sigma_1\sigma_2\sigma_3\sigma_4$}{$\lambda_1\lambda_2$}。SPBDD 內部把生成元化簡成下面這組（與常見的 $XZZXI$ 等生成同一個群），$\sigma_i$ 依序對應它們；$\lambda=(\langle e,\bar Z\rangle,\langle e,\bar X\rangle)$，邏輯算符代表元也是程式選出的：
\[
g_1=XIXZZ,\quad g_2=ZIZYY,\quad g_3=IXZZX,\quad g_4=IZYYZ,\qquad \bar X=IIXYX,\quad \bar Z=IIZXZ .
\]
（碼距離 3，$\bar X,\bar Z$ 都是權重 3。）下表是十個單量子位元 Pauli 的標籤，它們就是前面的 $w_v$：

\begin{center}
\small
\begin{tabular}{@{}lllll@{}}
\toprule
$X_1$ & $Z_1$ & $X_2$ & $Z_2$ & $X_3$\\
\tg{0100}{00} & \tg{1000}{00} & \tg{0001}{00} & \tg{0010}{00} & \tg{0111}{10}\\
\midrule
$Z_3$ & $X_4$ & $Z_4$ & $X_5$ & $Z_5$\\
\tg{1001}{01} & \tg{1111}{01} & \tg{0101}{11} & \tg{1101}{10} & \tg{0110}{01}\\
\bottomrule
\end{tabular}
\end{center}

關鍵性質：\textbf{標籤是線性的。}一個錯誤的標籤，就是它所含各個單量子位元 Pauli 的標籤逐位元 XOR（$Y=XZ$，所以 $Y_4$ 的標籤是 $X_4$ 與 $Z_4$ 的 XOR）。例如 $Y_4X_5$：
\[
\tg{1111}{01}\oplus\tg{0101}{11}\oplus\tg{1101}{10}=\tg{0111}{00}.
\]

\subsection{錯誤集合與顯式做法}

取大小為 5 的集合
\[
\E=\{\,IIIII,\ XIIII,\ IZIII,\ IIXII,\ IIIYX\,\}.
\]
逐一算標籤（顯式做法）：

\begin{center}
\small
\begin{tabular}{@{}lll@{}}
\toprule
錯誤 & 標籤 \tg{$\sigma$}{$\lambda$} & 算法\\
\midrule
$IIIII$ & \tg{0000}{00} & 全為 0\\
$XIIII$ & \tg{0100}{00} & $X_1$\\
$IZIII$ & \tg{0010}{00} & $Z_2$\\
$IIXII$ & \tg{0111}{10} & $X_3$\\
$IIIYX$ & \tg{0111}{00} & $Y_4X_5$（如上）\\
\bottomrule
\end{tabular}
\end{center}

最後兩列的症狀都是 \texttt{0111}，邏輯特徵卻不同（\texttt{10} 對 \texttt{00}）。驗證：$IIXII\cdot IIIYX=IIXYX=\bar X$，的確是一個邏輯算符。所以\textbf{答案是「有不等價對」}，這一對就是 $(IIXII,\ IIIYX)$。函式庫的 DP 與 compose 兩種方法都給出同樣的結論與同一對見證。

五個元素時顯式做法當然最簡單。問題在於真實的錯誤集合（例如所有權重 $\le 4$ 的錯誤）有 $10^8$--$10^9$ 個元素，無法逐一列出；決策圖只用幾百到幾千個節點就能「壓縮地」存下整個集合。DP 的目的就是\emph{不展開集合，直接在壓縮表示上算出標籤集合}。

\subsection{決策圖長什麼樣}

決策圖（BDD）把集合表示成一張「問答圖」：依固定順序（$X_1,Z_1,X_2,Z_2,\dots$）逐一問每個變數「這個單量子位元 Pauli 有沒有出現在錯誤裡？」虛線表示答「沒有」（取 0），實線表示答「有」（取 1）。從根節點走到 $\top$ 的每一條路徑就是集合中的一個元素。（走到 $\bot$ 的路徑不屬於集合。）

我們的 $\E$ 在真正的 BDD 裡有 19 個節點，但其中很多節點只是「這個變數必須為 0」的傳遞。把它們省略（它們不改變下面算出的標籤）後，只剩 6 個真正有分岔的節點，圖~\ref{fig:bdd} 畫的就是它們；每條路徑對應一個元素：

\begin{center}
\small
\begin{tabular}{@{}ll@{}}
\toprule
路徑 & 元素\\
\midrule
一路選虛線 & $IIIII$\\
在 $X_1$ 選實線，其餘為 0 & $XIIII$\\
在 $Z_2$ 選實線 & $IZIII$\\
在 $X_3$ 選實線 & $IIXII$\\
在 $X_4$、$Z_4$、$X_5$ 都選實線 & $IIIYX$\\
\bottomrule
\end{tabular}
\end{center}

\begin{figure}[ht]
\centering
\begin{tikzpicture}[
  node distance=1.25cm,
  dec/.style={circle,draw,minimum size=8mm,inner sep=0pt},
  term/.style={rectangle,draw,rounded corners,minimum height=6.5mm,inner sep=3pt,font=\small},
  img/.style={anchor=east,font=\small},
  zero/.style={-{Stealth},dashed},
  one/.style={-{Stealth}}]
  \node[dec] (r) at (0,0)    {$X_1$};
  \node[dec] (a) at (0,-1.6) {$Z_2$};
  \node[dec] (b) at (0,-3.2) {$X_3$};
  \node[dec] (c) at (0,-4.8) {$X_4$};
  \node[dec] (d) at (0,-6.4) {$Z_4$};
  \node[dec] (e) at (0,-8.0) {$X_5$};
  \node[term] (t) at (0,-9.6) {$\top$：$\{\tg{0000}{00}\}$};
  \node[term] (t1) at (3.6,0)    {其餘全 0 $\to\top$};
  \node[term] (t2) at (3.6,-1.6) {其餘全 0 $\to\top$};
  \node[term] (t3) at (3.6,-3.2) {其餘全 0 $\to\top$};
  \node[term] (t4) at (3.6,-4.8) {其餘全 0 $\to\top$};
  \node[term] (f1) at (3.6,-6.4) {$\bot$};
  \node[term] (f2) at (3.6,-8.0) {$\bot$};
  \draw[zero] (r) -- (a); \draw[one] (r) -- (t1);
  \draw[zero] (a) -- (b); \draw[one] (a) -- (t2);
  \draw[zero] (b) -- (c); \draw[one] (b) -- (t3);
  \draw[zero] (c) -- (t4); \draw[one] (c) -- (d);
  \draw[zero] (d) -- (f1); \draw[one] (d) -- (e);
  \draw[zero] (e) -- (f2); \draw[one] (e) -- (t);
  \node[img] at (-0.7,0)    {$I=\{\tg{0000}{00},\tg{0100}{00},\tg{0010}{00},\tg{0111}{00},\tg{0111}{10}\}$};
  \node[img] at (-0.7,-1.6) {$I=\{\tg{0000}{00},\tg{0010}{00},\tg{0111}{00},\tg{0111}{10}\}$};
  \node[img] at (-0.7,-3.2) {$I=\{\tg{0000}{00},\tg{0111}{00},\tg{0111}{10}\}$};
  \node[img] at (-0.7,-4.8) {$I=\{\tg{0000}{00},\tg{0111}{00}\}$};
  \node[img] at (-0.7,-6.4) {$I=\{\tg{1000}{01}\}$};
  \node[img] at (-0.7,-8.0) {$I=\{\tg{1101}{10}\}$};
\end{tikzpicture}
\caption{$\E$ 的（省略「必須為 0」節點後的）決策圖。虛線＝取 0，實線＝取 1。每個節點左邊的 $I$ 是「從這個節點往下走，能累積出的所有標籤」，由下往上依序算出（見下一小節）。}\label{fig:bdd}
\end{figure}

\subsection{DP：由下往上，每個節點只做一件事}

對每個節點 $u$，我們想知道它的\textbf{後綴像} $I(u)$：「從 $u$ 往下走到 $\top$ 的所有路徑，沿途經過的實線所對應的單量子位元 Pauli，其標籤 XOR 起來，可能是哪些值？」這就是 $u$ 以下所有元素的標籤集合。計算只有一條規則：

\begin{center}
\fbox{\parbox{0.9\textwidth}{\centering
$I(u)\;=\;\underbrace{I(\text{虛線的子節點})}_{\text{這個變數沒選：標籤不變}}\ \cup\ \underbrace{\bigl(w_v\oplus I(\text{實線的子節點})\bigr)}_{\text{這個變數選了：每個標籤都再 XOR 上 }w_v}$}}
\end{center}
其中 $w_v$ 是表中該變數的標籤；$\bot$ 的像是空集合，$\top$ 的像是 $\{\tg{0000}{00}\}$（什麼都沒選，標籤為零）。

照著圖由下往上算：
\begin{enumerate}[leftmargin=2em]
\item $X_5$ 節點：虛線通向 $\bot$（空集合），實線通向 $\top$。$I=\varnothing\cup(\tg{1101}{10}\oplus\{\tg{0000}{00}\})=\{\tg{1101}{10}\}$。
\item $Z_4$ 節點：虛線 $\bot$；實線通向上一步。$I=\tg{0101}{11}\oplus\{\tg{1101}{10}\}=\{\tg{1000}{01}\}$。
\item $X_4$ 節點：虛線通向「其餘全 0」，像是 $\{\tg{0000}{00}\}$；實線通向上一步。$I=\{\tg{0000}{00}\}\cup\bigl(\tg{1111}{01}\oplus\tg{1000}{01}\bigr)=\{\tg{0000}{00},\tg{0111}{00}\}$。（這正是 $IIIII$ 與 $IIIYX$ 兩個後綴的標籤。）
\item $X_3$ 節點：虛線通向上一步，實線通向「其餘全 0」。$I=\{\tg{0000}{00},\tg{0111}{00}\}\cup\{\tg{0111}{10}\}$。\textbf{這裡兩個標籤 \tg{0111}{00} 與 \tg{0111}{10} 第一次同時出現：症狀相同、邏輯特徵不同。}
\item $Z_2$ 節點：$I=(\text{上一步})\cup\{\tg{0010}{00}\}$。
\item $X_1$ 節點（根）：$I=(\text{上一步})\cup\{\tg{0100}{00}\}$，也就是五個元素的完整標籤集合 $G$。
\end{enumerate}
這個流程與表中直接算的結果完全一致，並且與函式庫對這張真實的 19 節點決策圖逐節點算出的像一致。

\subsection{$I$ 實際上怎麼存？為什麼強調「位移只是重新標號」}\label{sec:store}

\paragraph{$I$ 存成一個新的 BDD。}是的。圖~\ref{fig:bdd} 每個節點旁邊寫的 $I$ 看起來像「一堆標籤的清單」，實際上程式\emph{不}存清單：$I(u)$ 是標籤的 6 個位元 $(y_1,\dots,y_6)=(\sigma_1,\dots,\sigma_4,\lambda_1,\lambda_2)$ 上的一個布林函數「$y\in I(u)$ 嗎？」，以一個\emph{新的、獨立的 BDD} 存在同一個 BuDDy 管理器裡（程式中就是 \code{b}、\code{c}、\code{f} 這幾組變數）。輸入集合 $\E$ 的 BDD 在整個過程中只被讀取、不被修改；DP 每處理一個節點，就產生一個新的標籤 BDD，並用一個 \code{Bdd} 把手（引用計數的指標）記在該節點名下。例如節點 $X_5$ 的 $I=\{\tg{1101}{10}\}$ 只有一個點，它的 BDD 就是 6 個節點的一條鏈（每個變數一個節點，答案固定）。真實的大例子中 $I$ 可以是數億個標籤的集合，但只要結構規則，它的 BDD 只有幾千到幾十萬個節點——這就是不展開、直接在壓縮表示上算的意思。

\paragraph{DP 對 $I$ 做的兩件事。}由 $I(u)=I(\text{虛線子節點})\cup\bigl(w_v\oplus I(\text{實線子節點})\bigr)$，每個節點對標籤 BDD 做一次「位移」與一次「聯集」：
\begin{itemize}[leftmargin=1.5em]
\item \textbf{聯集}是普通的 BDD 運算（apply 的 OR）。它的結果可能比兩個輸入都大（最壞約為兩者節點數的乘積），所以\emph{這是 DP 中唯一可能讓圖變大的步驟}。
\item \textbf{位移} $w\oplus F$ 把集合中\emph{每一個}標籤的某幾個位元翻轉（$w$ 在哪些位置是 1，就翻轉哪些位置）。直覺上這會改變整個集合；但在 BDD 上它只是：\textbf{把標籤為被翻轉變數的節點，兩個子節點對調}（圖~\ref{fig:shift}）。節點數完全不變，只需掃一遍圖。
\end{itemize}

\begin{figure}[ht]
\centering
\begin{tikzpicture}[
  dec/.style={circle,draw,minimum size=7mm,inner sep=0pt,font=\small},
  term/.style={rectangle,draw,minimum size=5mm,inner sep=1pt,font=\footnotesize},
  zero/.style={-{Stealth},dashed},
  one/.style={-{Stealth}}]
  % left: F = {000, 011}
  \begin{scope}
    \node[dec] (y1) at (0,0) {$y_1$};
    \node[dec] (y2) at (0,-1.4) {$y_2$};
    \node[dec] (a) at (-1.4,-2.8) {$y_3$};
    \node[dec] (b) at (1.4,-2.8) {$y_3$};
    \node[term] (f0) at (1.2,0) {$\bot$};
    \node[term] (at) at (-2.0,-4.1) {$\top$};
    \node[term] (af) at (-0.8,-4.1) {$\bot$};
    \node[term] (bf) at (0.8,-4.1) {$\bot$};
    \node[term] (bt) at (2.0,-4.1) {$\top$};
    \draw[zero] (y1)--(y2); \draw[one] (y1)--(f0);
    \draw[zero] (y2)--(a);  \draw[one] (y2)--(b);
    \draw[zero] (a)--(at);  \draw[one] (a)--(af);
    \draw[zero] (b)--(bf);  \draw[one] (b)--(bt);
    \node at (0,-5.0) {$F=\{000,\ 011\}$};
  \end{scope}
  \draw[-{Stealth},thick] (3.1,-2.0) -- node[above,font=\small]{翻轉 $y_2$} node[below,font=\small]{$w=010$} (4.9,-2.0);
  % right: w xor F = {010, 001}: children of the y2 node swapped
  \begin{scope}[xshift=8cm]
    \node[dec] (y1) at (0,0) {$y_1$};
    \node[dec] (y2) at (0,-1.4) {$y_2$};
    \node[dec] (a) at (-1.4,-2.8) {$y_3$};
    \node[dec] (b) at (1.4,-2.8) {$y_3$};
    \node[term] (f0) at (1.2,0) {$\bot$};
    \node[term] (at) at (-2.0,-4.1) {$\top$};
    \node[term] (af) at (-0.8,-4.1) {$\bot$};
    \node[term] (bf) at (0.8,-4.1) {$\bot$};
    \node[term] (bt) at (2.0,-4.1) {$\top$};
    \draw[zero] (y1)--(y2); \draw[one] (y1)--(f0);
    \draw[zero] (y2)--(b);  \draw[one] (y2)--(a);
    \draw[zero] (a)--(at);  \draw[one] (a)--(af);
    \draw[zero] (b)--(bf);  \draw[one] (b)--(bt);
    \node at (0,-5.0) {$w\oplus F=\{010,\ 001\}$};
  \end{scope}
\end{tikzpicture}
\caption{位移 $w\oplus F$ 在 BDD 上的樣子（虛線＝取 0，實線＝取 1）：翻轉 $y_2$ 就是把 $y_2$ 節點的兩個子節點對調，其餘節點與邊完全不動，節點數不變。}\label{fig:shift}
\end{figure}

\paragraph{為什麼要特別強調這一點。}因為這正是 DP 能避開 compose 方法的瓶頸之處。
\begin{itemize}[leftmargin=1.5em]
\item compose 方法要把 BDD 的每個變數代換成\emph{其他變數的 XOR}（例如 $y_i:=x_i\oplus y_i$）。在 BDD 上這可能讓圖大指數倍（第~\ref{sec:compose} 節的例子：$n+2$ 個節點的函數變成至少 $2^n$ 個節點）。代價取決於「代換後的結果有多大」，事先無法控制。
\item DP 把同一件事換了一個角度：不去改寫輸入集合的變數，而是讓「錯誤的標籤是各單量子位元標籤的 XOR」這個線性性，變成對標籤集合的\emph{常數位移}。位移在 BDD 上是對調子節點，\emph{保證不增加節點數}。
\item 於是 DP 中每個節點只剩一個可能造成成長的運算——聯集，而且聯集的輸入都是「某個後綴的標籤集合」，它們通常比 compose 在 $2n$ 個變數上產生的中間圖小得多。DP 的成本因此只取決於各節點的 $I$ 有多大（第~\ref{sec:complexity} 節）。
\end{itemize}

\paragraph{實作上的小註記。}程式用 \code{Bdd::compose} 把每個被翻轉的變數 $y_t$ 代換成它自己的否定 $\neg y_t$ 來做位移，這在數學上等同於對調子節點（結果大小一樣）；但 BuDDy 的向量代換（\code{bdd_veccompose}）實作上不保證是線性時間，所以文件只宣稱「位移的結果大小不變」，執行時間只宣稱是圖大小的多項式。

\subsection{最後一步：檢查「症狀相同、特徵不同」}

得到 $G=\{\tg{0000}{00},\tg{0100}{00},\tg{0010}{00},\tg{0111}{00},\tg{0111}{10}\}$ 後，對每個特徵位元 $\lambda_j$ 檢查：有沒有某個症狀 $\sigma$，既出現在「$\lambda_j=0$ 的標籤」中、也出現在「$\lambda_j=1$ 的標籤」中（引理~\ref{lem:bits}）。對 $\lambda_1$：$\lambda_1=0$ 的標籤其症狀為 \texttt{0000, 0100, 0010, 0111}；$\lambda_1=1$ 的標籤只有 \texttt{0111}。交集是 $\{\texttt{0111}\}\neq\varnothing$，故有不等價對。最後用症狀 \texttt{0111} 回到原集合取元素：$\E\cap\{\text{症狀為 }\texttt{0111}\}=\{IIXII,\,IIIYX\}$，就是見證。

若把 $IIIYX$ 換成 $IIIIX$（標籤 \tg{1101}{10}，症狀 \texttt{1101} 不與其他元素重複），則 $G$ 的五個症狀兩兩不同，檢查不成立，答案為「無不等價對」。

\subsection{這個例子想說明的}
\begin{itemize}[leftmargin=1.5em]
\item 「位移」$w_v\oplus I(\cdot)$ 只是把集合裡每個標籤的某幾個位元翻轉；在決策圖上，這是重新標號，不會讓圖變大（第~\ref{sec:shift} 節）。
\item 每個節點只算一次，不論有多少條路徑經過它。五個元素時看不出差別；當集合有 $10^9$ 個元素而圖只有一千個節點時，DP 做的就只是約一千次「聯集＋位移」。
\item 這張圖上每個節點的 $I$ 都很小；在大集合上，$I$ 是否仍然小（結構是否保持）決定 DP 的成本，這是第~\ref{sec:complexity} 節的複雜度分析與第~\ref{sec:meas} 節量測在討論的事。
\end{itemize}

%% =====================================================================
\section{Compose 方法（參考實作）}\label{sec:compose}

把 $\{g_i,\bar X_j,\bar Z_j\}$ 補成辛基底（加入去穩定子 $d_i$），令 $T$ 為可逆的 $2n\times2n$ 矩陣，把 $e$ 送到座標 $(a\mid b\mid c\mid f)$，其中 $b=\sigma(e)$、$(c,f)=\lambda(e)$，$a$ 為與去穩定子配對的係數。$T(\E)$ 的特徵函數為 $\chi_\E\circ T^{-1}$，所以要對 $\chi_\E$ 的\emph{全部 $2n$ 個變數同時}代換：把變數 $y_v$ 代換成「$T^{-1}$ 第 $v$ 列所選出的新變數」的 XOR（\code{to_code_coordinates}），最後對 $a$ 區塊做存在量化（\code{relation_by_compose}）。

這個代換的代價取決於\emph{結果}的大小，而結果大小與輸入大小無關。例：$f(x,y)=\bigwedge_i\neg y_i$ 只有 $n+2$ 個節點，但代換 $y_i:=x_i\oplus y_i$ 會把它變成相等函數 $\bigwedge_i(x_i\leftrightarrow y_i)$，在變數順序 $x_1..x_n,y_1..y_n$ 下至少需要 $2^n$ 個節點。基底變換的列若含很多個 1，就會在 $2n$ 個變數的尺度上發生這種情形，而穩定子區塊到最後才被丟棄。

%% =====================================================================
\section{DP 方法}\label{sec:dp}

\subsection{想法}

只需要\emph{像}（image）$\pi(\E)$，而 $\pi$ 是線性的。把 $e=\sum_{v=0}^{2n-1}e_v u_v$ 寫成單位向量 $u_v$ 的和（$v=2q$ 對應單量子位元 Pauli $X_q$，$v=2q+1$ 對應 $Z_q$），並令
\[
w_v=\pi(u_v)\in\F_2^{m}\quad(\text{變數 }v\text{ 的「翻轉向量」}),\qquad
\pi(e)=\bigoplus_{v}e_v\,w_v .
\]
$\E$ 的 BDD 依層序（level order）對各變數做 Shannon 展開來列舉賦值，所以一個節點的像可以由其兩個子節點的像\emph{加上一個常數向量}得到，完全不必把某個變數改寫成其他變數的 XOR。

\subsection{定義}

設管理器的層為 $0,\dots,N-1$（$N=2n$），$\var(\ell)$ 為第 $\ell$ 層的變數。對 $\F_2^m$ 上的布林函數 $F$ 與向量 $w$，定義\emph{位移}（shift）
\[
(w\oplus F)(y)=F(y\oplus w),\qquad\text{即}\qquad w\oplus F=\{\,w\oplus t:t\in F\,\}.
\]
\paragraph{決策圖的記號。}BDD 有兩個\emph{終端節點}（terminal）：$\top$（讀作 top，「真」終端）與 $\bot$（讀作 bottom，「假」終端）。從根節點沿著各變數的取值往下走，若走到 $\top$，表示這組賦值\emph{屬於}集合 $\E$；若走到 $\bot$，表示\emph{不屬於}。所以 $\bot$ 代表空集合 $\emptyset$（沒有任何賦值被接受），$\top$ 則代表「剩下的變數取什麼值都被接受」。對非終端節點 $u$，$\mathrm{top}(u)$ 是 $u$ 所測試的變數（注意這個 $\mathrm{top}$ 是函數名稱，與終端 $\top$ 無關），$\low(u)$ 與 $\high(u)$ 分別是該變數取 $0$ 與取 $1$ 時所走向的子節點。$\oplus$ 表示 $\F_2$ 上的加法（位元 XOR）。

對 $\E$ 的 BDD 中的節點 $u$ 與層 $\ell\le N$，令 $\E_{u,\ell}$ 為：$u$ 所接受的、對第 $\ell,\dots,N-1$ 層變數的賦值之集合（$\ell$ 層以上的變數被忽略；被圖跳過的變數視為不受限制）。\emph{後綴像}（suffix image）定義為
\[
I(u,\ell)=\Bigl\{\bigoplus_{\ell'\ge\ell}e_{\var(\ell')}\,w_{\var(\ell')}\ :\ e\in\E_{u,\ell}\Bigr\}\subseteq\F_2^{m}.
\]
則 $G(\E)=I(\ra,0)$。

\subsection{遞迴式}\label{sec:rec}

令 $v=\var(\ell)$：
\begin{center}
\begin{tabular}{@{}p{0.45\textwidth}l@{}}
\toprule
情形 & $I(u,\ell)$\\
\midrule
$u=\bot$ & $\emptyset$\\
$\ell=N$（$u=\top$） & $\{0\}$\\
$u$ 是位於第 $\ell$ 層的節點 & $I(\low(u),\ell{+}1)\ \cup\ \bigl(w_v\oplus I(\high(u),\ell{+}1)\bigr)$\\
其他（第 $\ell$ 層被跳過，或 $\ell<N$ 時 $u=\top$，後者表示剩下的變數全部不受限制） & $I(u,\ell{+}1)\ \cup\ \bigl(w_v\oplus I(u,\ell{+}1)\bigr)$\\
\bottomrule
\end{tabular}
\end{center}

\begin{theorem}[正確性]\label{thm:correct}
遞迴式計算出的就是 3.2 節所定義的 $I(u,\ell)$。特別地，$I(\ra,0)=\pi(\E)=G(\E)$。
\end{theorem}
\begin{proof}
對 $N-\ell$ 做歸納。$\ell=N$ 時唯一的後綴賦值是空賦值，若 $u$ 接受它（$u=\top$）則像為 $\{0\}$，否則為 $\emptyset$。$\ell<N$ 時，依第 $\ell$ 層變數的取值 $b$ 切分 $\E_{u,\ell}$。若 $u$ 測試該變數，則 $b=0$ 的賦值恰為 $\E_{\low(u),\ell+1}$，貢獻的像是 $I(\low(u),\ell+1)$；$b=1$ 的賦值恰為 $\E_{\high(u),\ell+1}$，由 $\bigoplus$ 的線性性，其貢獻為這些像各自加上 $w_v$，即 $w_v\oplus(\cdot)$。若 $u$ 不測試該變數，則 $b$ 取兩值時剩下的限制條件同為 $\E_{u,\ell+1}$，得到最後一列。
\end{proof}

\subsection{位移只是重新標號}\label{sec:shift}

$w\oplus F$ 可由 $F$ 的決策圖得到：對每個以 $\supp(w)$ 中變數為標籤的節點，交換其兩個子節點。這是節點上的一個雙射，且保持化簡性（reducedness）——兩節點相等若且唯若其像相等，節點冗餘若且唯若其像冗餘——因此
\[
\lvert w\oplus F\rvert=\lvert F\rvert\quad\text{（節點數完全相同）}.
\]
實作上以「把每個被翻轉的變數代換成它自己的否定」（\code{Bdd::compose}，代換為文字 $\neg y_t$）來表達，絕不使用「多個變數的 XOR」。

\subsection{演算法}\label{sec:algo}

遞迴式本身（對 $(u,\ell)$ 定義）：
\begin{algorithm}[H]
\caption{\textsc{Image}$(u,\ell)$：後綴像}
\begin{algorithmic}[1]
\If{$u=\bot$} \State \Return $\emptyset$ \EndIf
\If{$\ell=N$} \State \Return $\{0\}$ \Comment{$m$ 個目標變數上的方塊 $y=0$} \EndIf
\State $v\gets\var(\ell)$
\If{$u$ 是節點且 $\lev(\mathrm{top}(u))=\ell$}
  \State \Return \textsc{Image}$(\low(u),\ell{+}1)\ \cup\ \textsc{Shift}(\textsc{Image}(\high(u),\ell{+}1),w_v)$
\Else \Comment{第 $\ell$ 層被跳過}
  \State $H\gets$ \textsc{Image}$(u,\ell{+}1)$
  \State \Return $H\cup\textsc{Shift}(H,w_v)$ \Comment{$w_v=0$ 時即為 $H$}
\EndIf
\end{algorithmic}
\end{algorithm}

函式庫實際採用\emph{逐層}（level by level）求值：

\begin{algorithm}[H]
\caption{逐層求值（函式庫所用）}
\begin{algorithmic}[1]
\State $\mathrm{need}[0]\gets\{\ra\}$
\For{$\ell=0,\dots,N-1$} \Comment{向下掃描：哪些 $(u,\ell)$ 會被用到}
  \State 把 $\mathrm{need}[\ell]$ 中每個 $u$ 的子節點（若第 $\ell$ 層被跳過則為 $u$ 本身）加入 $\mathrm{need}[\ell{+}1]$，略去 $\bot$
\EndFor
\State $\mathit{next}\gets\{(u,N)\mapsto\{0\}: u\in\mathrm{need}[N]\}$
\For{$\ell=N-1$ \textbf{downto} $0$} \Comment{向上掃描}
  \State $\mathit{cur}\gets\{(u,\ell)\mapsto\textsc{Image}(u,\ell)\text{ 由 }\mathit{next}\text{ 算出}: u\in\mathrm{need}[\ell]\}$
  \State $\mathit{next}\gets\mathit{cur}$ \Comment{第 $\ell{+}1$ 層的像於此釋放}
\EndFor
\State $G\gets\mathit{next}[\ra]$
\end{algorithmic}
\end{algorithm}

\begin{algorithm}[H]
\caption{\code{find_inequivalent_pair}$(\E)$}
\begin{algorithmic}[1]
\State $G\gets\textsc{Image}(\ra(\E),0)$ \Comment{建在 $b,c,f$ 變數上}
\State 對 $G$ 執行引理~\ref{lem:bits} 的 $2k$ 個單一位元檢驗
\If{沒有任何檢驗成立} \State \Return 「無不等價對」\EndIf
\State 取出有問題的症狀 $s$，以及纖維 $G\wedge(\sigma=s)$ 中兩個特徵 $\lambda_1\ne\lambda_2$
\State \Return $\E\cap\code{with_syndrome}(s)\cap\code{with_logical_signature}(\lambda_i)$ 中的任意元素（$i=1,2$）
\end{algorithmic}
\end{algorithm}

實作細節（\code{src/stabilizercode.cpp}）：
\begin{itemize}[leftmargin=1.5em]
\item \textbf{預先計算。}建構子對 $2n$ 個變數各存一份清單 \code{flips_}：列出 $w_v$ 為 1 的目標變數（\code{b_var}、\code{c_var}、\code{f_var}）。共 $2n\cdot m$ 次辛內積，$O(n^3)$ 位元運算，不超過建構子原本就要做的線性代數。去穩定子區塊不必儲存：它只影響 $a$ 座標，而 $a$ 在對 $S$ 取商時被丟棄。
\item \textbf{目標變數}與 compose 路徑使用相同的 \code{b_var/c_var/f_var} 編號，所以兩條路徑回傳\emph{完全相同的正規（canonical）決策圖} $G$（已有測試），multiplicity 檢驗與見證（witness）的取出也由兩者共用。
\item \textbf{見證}現在在原座標中讀出：\code{errors & with_syndrome(s)} \code{& with_logical_signature(lam)}，因為 $G$ 已忘記每個點是由哪個錯誤產生的。兩種方法皆如此。
\item \textbf{重排序（reordering）。}DP 在讀取 $\E$ 的各層的同時建構其他決策圖；中途若動態重排序，層就會移動。因此 \code{ReorderPause} 在呼叫期間關閉動態重排序、結束後還原（為此新增了 \code{Manager::dynamic_}\allowbreak\code{reordering()}）。
\item \textbf{求值順序。}遞迴式以逐層由下往上的方式計算，第 $\ell$ 層的像一算好，第 $\ell+1$ 層的像就立即釋放。深度優先加 memo 的遞迴會算出相同的圖，但會把每個 $(u,\ell)$ 的像一直留到結束；其代價見第~\ref{sec:why} 節（在最大的權重球上約為 10 倍）。
\end{itemize}

%% =====================================================================
\section{複雜度}\label{sec:complexity}

記號：$\lvert f\rvert$ 為輸入決策圖的節點數，$N=2n$，$m=r+2k$，
\[
\skipc(f)=\sum_{(u\to c)}\bigl(\lev(c)-\lev(u)-1\bigr)
\]
（對 $f$ 的所有邊求和，即被跳過的層數）。

\begin{itemize}[leftmargin=1.5em]
\item \textbf{$(u,\ell)$ 配對數。}每個被用到的 $(u,\ell)$ 若不是 $(u,\lev(u))$，就是落在「由某條通向 $u$ 的邊進入」的一串被跳過的層上。因此
\[
M\le\lvert f\rvert+\skipc(f)\le\lvert f\rvert\,(1+N),
\]
對很少跳層的圖，$M\approx\lvert f\rvert$。終端 $\top$ 至多貢獻 $N$ 個配對。
\item \textbf{每個配對的工作量。}至多一次位移與一次聯集（union）。位移的結果與原圖節點數完全相同（第~\ref{sec:shift} 節）；就重新標號而言，理想情況下每個被翻轉的變數只需一趟掃描，但實作所用的 BuDDy \code{bdd_veccompose} 不保證線性，所以只宣稱「圖大小的多項式」。$A$ 與 $B$ 的聯集代價為 $O(\lvert A\rvert\,\lvert B\rvert)$，結果至多這麼多節點。
\item \textbf{總計。}
\[
T_{\mathrm{DP}}=\sum_{(u,\ell)}O\Bigl(\lvert I(\mathrm{lo})\rvert\cdot\lvert I(\mathrm{hi})\rvert+\lvert w_v\rvert\,\lvert I(\mathrm{hi})\rvert\Bigr),
\qquad
S_{\mathrm{DP}}=O\Bigl(\max_\ell\sum_{u\text{ 位於第 }\ell\text{ 層}}\lvert I(u,\ell)\rvert\Bigr).
\]
（逐層求值只需同時保存相鄰兩層；深度優先的 memo 版本則需對\emph{所有} $(u,\ell)$ 求和。）每個 $I(u,\ell)$ 都是 $\F_2^m$ 的子集，故每張圖至多 $2^m$ 個點，且 $I(\ra,0)=G$。此界對後綴像是\emph{輸出敏感}的：只要 $\E$ 各後綴的像在 $(\sigma,\lambda)$ 空間中仍保持結構化，它就小；一般情形沒有先驗的多項式界——本文並不宣稱此判定問題一般而言是易解的。
\item \textbf{Multiplicity 檢驗}（共用）：$2k$ 輪，每輪兩次合取、兩次對 $2k$ 個變數的量化，各為 $O(\lvert G\rvert)$（不計 apply 的代價）。在多數量測中相對於建構 $G$ 都很小；$d{=}11,t{=}4$ 的權重球上約 3.5 秒，已占總時間的約三分之一，但仍遠小於 compose 的代換。
\item \textbf{Compose 路徑}（比較用）：對 $N$ 個變數做一次同時代換，每個代入的 XOR 長度總和等於 $T^{-1}$ 中 1 的個數（至多 $N^2$），之後再做 $\exists a$。其執行時間是被代換後、位於 $2n$ 個變數上的決策圖大小的函數，可以比 $\lvert f\rvert$ 與 $\lvert G\rvert$ 大指數倍（第~\ref{sec:compose} 節）。
\end{itemize}

DP 在實務上獲勝的一句話理由：兩種方法最終都得到 $r+2k=n+k$ 個變數上的同一個 $G$，但 compose 必須經過一張由 XOR 耦合的、$2n$ 個變數上的決策圖，而 DP 只持有 $n+k$ 個目標變數上的圖，並且只以重新標號在它們之間移動。

%% =====================================================================
\section{與逐位置建構的關係}

若錯誤集合本身就是由小集合的和建出來的——「至多 $t$ 個位置，每個位置貢獻少數幾種 Pauli 之一」——則像也可以用分層集合建構：
\[
G_j\leftarrow G_j\cup\bigcup_p\bigl(\pi(p)\oplus G_{j-1}\bigr),
\]
完全不需要輸入決策圖。在權重球上，該建構與 DP 得到的 $G$ 相同（$G$ 是正規的，節點數一致）。DP 不需要知道 $\E$ 是怎麼建出來的，因此適用於\emph{任意}的 \code{PauliSet}。

%% =====================================================================
\section{API 變更}

\begin{lstlisting}[language=C++]
enum class StabilizerCode::PairMethod { Dp, Compose };
bool has_inequivalent_pair (const PauliSet&, PairMethod = PairMethod::Dp) const;
std::optional<Pair> find_inequivalent_pair(const PauliSet&,
                                           PairMethod = PairMethod::Dp) const;
bool Manager::dynamic_reordering() const;     // 新增，唯讀
\end{lstlisting}

compose 的實作保留為 \code{PairMethod::Compose}，作為參考實作與基準測試之用。

\section{測試}

\code{test/stabilizercode_test.cpp} 的「the DP and the compose route agree」一節：對位元翻轉碼、五量子位元碼與 Steane 碼各取 60 個隨機陪集（coset）聯集（安全與不安全的集合都會出現），比較兩種方法的判定結果，並獨立驗證兩者回傳的每個見證（屬於該集合、症狀相同、特徵不同、乘積落在 $N(S)\setminus S$）；已知答案的權重球；開啟動態重排序時的行為。所有既有測試（對所有子集的暴力驗證、乘積式 $E{*}E\cap(N(S)\setminus S)$）不經修改全部通過。

%% =====================================================================
\section{量測}\label{sec:meas}

所有執行：同一個雲端容器（2 vCPU、7~GB）、\code{g++ -O2}、BuDDy 節點表 $2^{22}$ 個節點／快取 $2^{20}$、預設變數順序、不做重排序，每列只跑一次（無重複）。DP 的數字是第~\ref{sec:algo} 節逐層版本、單獨在機器上執行的結果；compose 的數字是較早量到的（部分執行時另有一個工作同時在跑）。「compose」指 \code{PairMethod::Compose}，「DP」指 \code{PairMethod::Dp}；兩者都跑完的列，判定結果一律相同。有時兩個核心上同時各跑一個工作，會讓兩者都變慢，因此時間請以約兩倍的誤差來讀。第~\ref{sec:vsexp} 節的數據取自較晚的另一個量測時段，該時段機器明顯較慢（同一個 DP 在 $d{=}11,t{=}4$ 的球上，先前 11.6 s、該時段 18.9 s），所以\emph{不同時段的表之間不要直接比較絕對時間}，只在同一張表內比較。節點表大小的影響同樣重要：在 BuDDy 較小的預設表（$2^{20}$）上，compose 在 $d{=}7,t{=}3$ 600 秒內跑不完，但在 $2^{22}$ 上只需 111 秒。

\paragraph{(a) 權重球。}旋轉表面碼上所有權重 $\le t$ 的 Pauli（答案：無不等價對，因為 $d>2t$）。
\begin{center}
\begin{tabular}{@{}rrrrrr@{}}
\toprule
$d$ & $t$ & $\lvert\E\rvert$ & 圖節點數 & DP & compose\\
\midrule
5 & 2 & 2,776 & 139 & 0.003 s & 0.026 s\\
7 & 2 & $1.07\times10^4$ & 283 & 0.016 s & 1.55 s\\
7 & 3 & $5.08\times10^5$ & 369 & 0.082 s & 111.1 s\\
9 & 3 & $2.33\times10^6$ & 625 & 0.37 s & 600 s 內未完成\\
9 & 4 & $1.37\times10^8$ & 771 & 2.5 s & 未執行\\
11 & 4 & $6.96\times10^8$ & 1,171 & 11.6 s & 未執行\\
13 & 4 & $2.7\times10^9$ & 未量測 & 106 s & 未執行\\
\bottomrule
\end{tabular}
\end{center}

\paragraph{(b) 無結構集合。}同一個碼上，取 $K$ 個隨機陪集的聯集，每個陪集有 $G$ 個權重 $\le W$ 的隨機生成元（\code{pair_benchmark cosets}；固定隨機種子）。四列的答案皆為無不等價對。
\begin{center}
\begin{tabular}{@{}rrrrrr@{}}
\toprule
$d$ & $K,G,W$ & $\lvert\E\rvert$ & 圖節點數 & DP & compose\\
\midrule
5 & 20, 6, 4 & 1,280 & 7,265 & 0.017 s & 0.289 s\\
7 & 50, 8, 5 & $1.28\times10^4$ & 113,772 & 0.53 s & 65.2 s\\
9 & 100, 8, 5 & $2.56\times10^4$ & 387,759 & 2.1 s & 600 s 內未完成\\
11 & 300, 8, 5 & $7.67\times10^4$ & 1,569,997 & 20.5 s & 未執行\\
\bottomrule
\end{tabular}
\end{center}

\paragraph{(c) 電路的錯誤集合。}一輪帶旗標（flag）量子位元的症狀萃取中，所有故障（至多 $t$ 個故障位置；輸入資料錯誤也算故障）且限於旗標全為 0 的輸出，並把每個測量記錄以 Pauli 分量保留在被測量的量子位元上（因此交給 \code{StabilizerCode} 的碼是 $S\otimes I$ 再加上每個被測量位元一個單量子位元生成元，共 $N=n+\#\text{輔助}+\#\text{旗標}$ 個量子位元）。這些資料由包在函式庫外的獨立實驗程式產生，不在本儲存庫的程式碼中；故障集合以 \code{cx}／\code{fault_inject} 與 \code{|} 建構。時間僅計 \code{find_inequivalent_pair}。
\begin{center}
\begin{tabular}{@{}lrrrrrr@{}}
\toprule
電路 & $N$ & $t$ & 集合節點數 & DP & compose & 顯式雜湊\\
\midrule
Steane，帶旗標 & 19 & 2 & 7,405 & 0.011 s & 0.125 s & 0.002 s\\
Golay $[[23,1,7]]$，帶旗標 & 67 & 1 & 11,591 & 0.022 s & 550 s & 0.004 s\\
Surface-5，帶旗標 & 65 & 2 & 322,470 & 0.66 s & 56.3 s & 0.023 s\\
Golay，帶旗標 & 67 & 2 & 705,953 & 3.2 s & 未完成 & 0.087 s\\
\bottomrule
\end{tabular}
\end{center}

\paragraph{數字說明了什麼。}
\begin{itemize}[leftmargin=1.5em]
\item 凡是 compose 跑得完的集合，DP 都比較快，從約 2 倍（最小的集合）到約 25,000 倍（Golay，$t{=}1$：550 s 對 0.022 s）。$d{=}7,t{=}3$ 的球為 1,350 倍，帶旗標的 surface-5 集合為 85 倍。compose 600 秒內跑不完的集合，DP 幾秒內解出（$d{=}9$ 的球 0.4 s、$d{=}9$ 的陪集聯集 2 s），$d{=}13,t{=}4$ 的球（$2.7\times10^9$ 個元素）也只要 106 s。
\item 當集合小到可以逐一列出時，DP 並不會讓決策圖方法贏過直接對故障組合做雜湊（(c) 的「顯式雜湊」一欄）；在元素數大的集合上情形相反。什麼時候哪一邊較好，第~\ref{sec:vsexp} 節在同一輪量測中同時量了時間與峰值記憶體。
\item 表 (c) 的 DP 欄只計 \code{find_inequivalent_pair}。建構故障集合本身要大得多（例如實驗程式中 Golay、$t{=}2$ 要 225 s），所以在電路層級的分析中，檢驗並不是最主要的成本。Golay、$t{=}2$ 那一列先前寫成「$\sim4\times10^6$」個節點，那其實是未過濾的故障集合的節點數；真正被檢驗的集合（旗標全為 0）有 705,953 個。
\item 兩種方法的 multiplicity 檢驗完全相同，上述時間皆已包含；巨大的差異來自 $G$ 的取得方式（對 compose 而言，另外單獨計時的代換步驟在 $d{=}7,t{=}3$ 的球與 Golay $t{=}1$ 上占 99\% 以上的時間）。
\end{itemize}

\subsection{為什麼最初的 DP 版本比逐層建構慢}\label{sec:why}

最初的實作用雜湊表 memo 住每個 $(u,\ell)$ 的像，並採深度優先求值。在 $d{=}11,t{=}4$ 的球上需要 113 s（下面的剖析執行為 175 s），而直接由逐量子位元的像建構 $G$ 的分層建構（$G_j\leftarrow G_j\cup\bigcup_p(\pi(p)\oplus G_{j-1})$，只適用於具有這種形式的集合）約需 24 s。把三種求值順序放進同一個程式剖析，每次單獨執行，同一台機器、同一個球、同樣的位移／聯集程式碼（A = memo 深度優先，B = 分層逐量子位元，C = 逐層，即第~\ref{sec:algo} 節；GC = BuDDy 垃圾回收）：
\begin{center}
\small
\begin{tabular}{@{}llrrrrr@{}}
\toprule
 & 節點表 & 總時間 & 位移 & 聯集 & GC 次數／時間 & 使用中節點（含垃圾）\\
\midrule
A，memo 深度優先 & $2^{22}$ & 175.0 s & 38.5 s & 80.4 s & 223 ／ 102.1 s & 15.3 M\\
B，分層 & $2^{22}$ & 24.2 s & 5.9 s & 12.0 s & 35 ／ 5.0 s & 4.6 M\\
C，逐層 & $2^{22}$ & 13.0 s & 2.7 s & 3.9 s & 16 ／ 3.0 s & 4.6 M\\
A，memo 深度優先 & $2^{25}$ & 10.3 s & 2.9 s & 3.9 s & 0 ／ 0 s & 22.2 M\\
B，分層 & $2^{25}$ & 21.5 s & 6.9 s & 11.8 s & 2 ／ 0.7 s & 8.5 M\\
C，逐層 & $2^{25}$ & 10.7 s & 3.0 s & 4.2 s & 0 ／ 0 s & 22.2 M\\
\bottomrule
\end{tabular}
\end{center}
（「使用中節點」為結束時 BuDDy 節點表中的使用量，含尚未回收的垃圾。位移或聯集的時間含它所觸發的垃圾回收，所以 GC 欄與前兩欄重疊。總時間中其餘的約 3.5 s 是 multiplicity 檢驗。A 與 C 做了相同的 1,646 次位移與 1,646 次聯集；B 做了 1,452 次位移與 1,936 次聯集。）

由此可知：
\begin{itemize}[leftmargin=1.5em]
\item \textbf{原因是記憶體被保留，不是運算次數。}memo 深度優先把 1,646 個中間像一直留到結束；節點表長到 15 M 個節點在使用中，而最終的 $G$ 只有 2.4 M（只保留 $t+1$ 層的分層建構為 4.6 M）。節點表只有 $2^{22}$ 時，BuDDy 必須垃圾回收並擴表 223 次，占了 175 s 中的 102 s；每次回收都要標記所有存活節點。節點表開到 $2^{25}$（不需任何回收）時，同一份程式只要 10.3 s，與逐層版本相同，也比分層建構快。
\item \textbf{解法是求值順序。}逐層求值時 DP 只持有相鄰兩層的像（此次執行中同時至多 14 個），且與節點表大小無關：$2^{22}$ 時 13.0 s，$2^{25}$ 時 10.7 s。這就是函式庫中的版本。
\item \textbf{DP 每次運算的工作量比分層建構小}（位移 2.7--3.0 s 對 5.9--6.9 s，聯集 3.9--4.2 s 對 11.8--12.0 s），雖然運算次數相近。原因尚未進一步調查。
\item 最初版本的執行時間會有波動（函式庫建置版本為 113 s，剖析程式為 175 s），因為成本取決於 BuDDy 何時剛好進行回收；這是同一個原因的另一個症狀。
\end{itemize}

\subsection{DP 與顯式列舉：時間與峰值記憶體}\label{sec:vsexp}

這一節回答「在什麼場景下 DP 比顯式列舉好」。數據由 \code{examples/dp_vs_explicit.cpp} 產生，兩種方法在同一個程式、同一台機器、同一輪量測中依序單獨執行（單執行緒）。

\paragraph{顯式基線。}逐一列舉集合的每個元素，用預先算好的 $\pi(\text{單量子位元 Pauli})$ 做增量 XOR 得到 $(\sigma,\lambda)$（每個元素一次 XOR），並維護雜湊表 $\sigma\mapsto\{\text{已見的 }\lambda\}$；某個 $\sigma$ 出現兩個不同的 $\lambda$ 即停止。金鑰是完整的 $\sigma$（無雜湊碰撞），每筆 16--32 位元組（依 $r$ 而定）。雜湊表上限設為 2~GB；若表會超過上限，就把症狀空間分成 $P=2,4,8,\dots$ 份、每趟只處理一份，並重新列舉全部元素——以時間換記憶體。表中「顯式」的時間是最後成功那一種 $P$ 的時間（括號內為含失敗嘗試的總時間），記憶體是整個行程的駐留記憶體峰值（雜湊表擴容時新舊表並存，所以 2~GB 的上限對應約 3.1~GB 的峰值）。這個基線\emph{被告知集合是怎麼建出來的}（直接依配方列舉），這對它有利；DP 只拿到集合的決策圖。DP 的記憶體包含 BuDDy 預先配置的 $2^{22}$ 節點表（行程啟動時就是 228~MB），所以小集合上 DP 的記憶體不會低於這個值。

\paragraph{(E1) 權重球（旋轉表面碼，所有權重 $\le t$ 的 Pauli）。}
\begin{center}
\small
\begin{tabular}{@{}rrrrrrrr@{}}
\toprule
$d$ & $t$ & $\lvert\E\rvert$ & DP 時間 & DP 記憶體 & 顯式時間 & $P$ & 顯式記憶體\\
\midrule
9  & 3 & $2.33\times10^6$ & 0.48 s & 228 MB & 0.44 s & 1 & 101 MB\\
11 & 3 & $7.84\times10^6$ & 1.75 s & 228 MB & 1.75 s & 1 & 389 MB\\
13 & 3 & $2.15\times10^7$ & 5.4 s & 229 MB & 6.5 s & 1 & 1,158 MB\\
15 & 3 & $5.08\times10^7$ & 13.2 s & 230 MB & 19.8 s & 1 & 3,078 MB\\
9  & 4 & $1.37\times10^8$ & 3.7 s & 228 MB & 45.7 s (83.2 s) & 2 & 3,108 MB\\
11 & 4 & $6.96\times10^8$ & 18.9 s & 237 MB & 399 s (536 s) & 8 & 3,108 MB\\
13 & 4 & $2.68\times10^9$ & 214 s & 357 MB & \multicolumn{3}{l}{80 分鐘內未完成（$P=64$，被中止）}\\
11 & 5 & $4.9\times10^{10}$ & 900 s 內未完成 & --- & \multicolumn{3}{l}{顯式未執行（外推需數天以上）}\\
\bottomrule
\end{tabular}
\end{center}

\noindent（$d{=}13,t{=}4$ 的顯式：事先由相異症狀數估出 $P=64$ 才夠放進 2~GB，直接以 $P=64$ 執行，80 分鐘仍未完成而被中止；以 $d{=}11,t{=}4$ 的每元素每趟約 72~ns 外推約需 3 小時。$d{=}11,t{=}5$ 的外推依同一速率，並假設 $P$ 隨 $\lvert\E\rvert$ 增加。）

\paragraph{(E2) 整塊區域。}$d{=}9$ 的旋轉表面碼上，集合為「支撐在一個 $a\times b$ 區塊內的所有 Pauli」，共 $4^{ab}$ 個元素（答案皆為無不等價對）。
\begin{center}
\small
\begin{tabular}{@{}lrrrrr@{}}
\toprule
區塊 & $\lvert\E\rvert$ & 決策圖節點數 & DP 時間 & 顯式時間 & 顯式記憶體\\
\midrule
$3\times3$ & $2.6\times10^5$ & 145 & $<0.01$ s & 0.007 s & 6 MB\\
$4\times3$ & $1.7\times10^7$ & 139 & $<0.01$ s & 0.49 s & 6 MB\\
$4\times4$ & $4.3\times10^9$ & 131 & $<0.01$ s & 155.6 s & 53 MB\\
$5\times4$ & $1.1\times10^{12}$ & 123 & $<0.01$ s & 未執行（外推約 11 小時） & \\
$5\times5$ & $1.1\times10^{15}$ & 113 & $<0.01$ s & 不可行 & \\
$6\times6$ & $4.7\times10^{21}$ & 91 & 0.009 s & 不可行 & \\
$8\times8$ & $3.4\times10^{38}$ & 35 & 0.054 s & 不可行 & \\
\bottomrule
\end{tabular}
\end{center}
（DP 的記憶體在所有列都是啟動時的 228~MB。顯式的外推取 $4\times4$ 的 36~ns／元素。）

\paragraph{(E3) 區塊內的權重球。}同一個碼，集合為「支撐在 $a\times a$ 區塊內、權重 $\le t$ 的 Pauli」。
\begin{center}
\small
\begin{tabular}{@{}rrrrrr@{}}
\toprule
區塊 & $t$ & $\lvert\E\rvert$ & DP 時間 & 顯式時間 & 顯式記憶體\\
\midrule
$6\times6$ & 3 & $1.99\times10^5$ & 0.034 s & 0.018 s & 11 MB\\
$6\times6$ & 4 & $4.97\times10^6$ & 0.119 s & 1.27 s & 197 MB\\
$5\times5$ & 5 & $1.40\times10^7$ & 0.194 s & 3.77 s & 389 MB\\
$6\times6$ & 5 & $9.66\times10^7$ & 0.62 s & 32.1 s & 3,077 MB\\
$6\times6$ & 6 & $1.52\times10^9$ & 5.3 s & 857 s（$P=8$；含失敗嘗試 1,061 s） & 3,108 MB\\
\bottomrule
\end{tabular}
\end{center}
（DP 的記憶體同樣是 228~MB。）

\paragraph{(E4) 無結構集合（隨機陪集聯集）。}與第~\ref{sec:meas} 節的表 (b) 同一種集合：$K$ 個陪集，每個有 $G$ 個隨機生成元、權重 $\le W$（固定種子）。
\begin{center}
\small
\begin{tabular}{@{}rrrrrrr@{}}
\toprule
$d$ & $K,G,W$ & $\lvert\E\rvert$ & 圖節點數 & DP 時間／記憶體 & 顯式時間 & 顯式記憶體\\
\midrule
7  & 50, 8, 5 & $1.28\times10^4$ & 113,772 & 0.96 s／235 MB & 0.001 s & 6 MB\\
9  & 100, 8, 5 & $2.56\times10^4$ & 387,759 & 3.5 s／251 MB & 0.003 s & 6 MB\\
11 & 300, 8, 5 & $7.67\times10^4$ & 1,569,997 & 36.4 s／337 MB & 0.006 s & 8 MB\\
7  & 20, 16, 4 & $1.18\times10^6$ & 562,008 & 21.0 s／266 MB & 0.31 s & 53 MB\\
9  & 20, 20, 4 & $2.10\times10^7$ & 4,885,255 & 300 s 內未完成 & 7.4 s & 773 MB\\
9  & 50, 24, 4 & $\le8.4\times10^8$ & 建圖就未完成 & 300 s 內未完成 & 300 s 內未完成 & \\
\bottomrule
\end{tabular}
\end{center}

\paragraph{由這些數字得到的結論。}
\begin{itemize}[leftmargin=1.5em]
\item \textbf{顯式的代價與 $\lvert\E\rvert$ 成正比，DP 的代價與決策圖及其像的大小有關。}顯式的時間是 $\Theta(\lvert\E\rvert)$ 次雜湊操作乘上趟數 $P$；記憶體是相異症狀數乘上每筆大小，除以 $P$。在權重球上相異症狀數約為 $0.7$--$0.8\,\lvert\E\rvert$（$d{=}11,t{=}4$：$5.3\times10^8$ 個症狀，若不分趟約需 12~GB），所以顯式在 $\lvert\E\rvert$ 超過約 $10^8$ 後就必須分趟，時間再乘上 $P$。DP 的時間隨 $(d,t)$ 增加得慢得多，記憶體幾乎不變（表 (E1) 的 DP 記憶體在 228--357~MB 間）。
\item \textbf{DP 明顯較好的場景：元素數極大、而決策圖小且像仍有結構。}(E1) 中 $t=4$ 的球：時間快 12--21 倍（$d{=}9$：3.7 s 對 45.7 s；$d{=}11$：18.9 s 對 399 s），峰值記憶體少 13 倍以上（237~MB 對 3.1~GB）。(E3) 的區塊內權重球：$6\times6,t{=}6$ 快約 160 倍。(E2) 的整塊區域：DP 在 0.01 秒內處理 $10^{38}$ 個元素，顯式在 $4\times4$（$4.3\times10^9$）就要 156 秒，更大的完全不可行。這類集合在 QEC 中很常見：局部（patch）上的任意錯誤、固定權重以內的錯誤。
\item \textbf{記憶體上的優勢比時間上的優勢出現得更早。}$t=3$ 的球上兩者時間相當（0.44 s 對 0.48 s；$d{=}13$、$d{=}15$ 上 DP 快 1.2--1.5 倍），但顯式的記憶體已從 101~MB 長到 $d{=}15$ 的 3~GB，而 DP 仍是約 230~MB。換句話說，若不想（或不能）用分趟來換記憶體，DP 在較小的集合上就已經划算。
\item \textbf{顯式明顯較好的場景：集合能在秒級內列完，而其決策圖相對於 $\lvert\E\rvert$ 並不小。}(E4) 中的隨機陪集聯集，$\lvert\E\rvert\le10^5$ 時顯式快約 1,000--6,000 倍（0.006 s 對 36.4 s），記憶體也只要 6--8~MB；到 $\lvert\E\rvert=2\times10^7$，DP 的決策圖達 490 萬個節點，300 秒內仍未完成，顯式只要 7.4 秒。DP 的代價取決於決策圖與各後綴像的大小，當這兩者與 $\lvert\E\rvert$ 同量級時，DP 就失去優勢。
\item \textbf{兩者都失效的場景：}元素數 $\sim10^9$ 且沒有結構（表 (E4) 最後一列）：顯式的雜湊表在 $P=1,2,4,8$ 都超過 2~GB 的上限、300 秒內也沒完成，DP 連集合的決策圖都建不出來。
\item \textbf{公平性的註記。}(i) 顯式基線拿到的是建構配方；若只給它決策圖，必須自己從圖上列舉元素，代價會更高（未量測）。(ii) 顯式很容易平行化（各趟、各分區彼此獨立），這裡只量了單執行緒；DP 目前也是單執行緒。(iii) 表 (E2) 的集合是線性子空間，若利用線性代數（高斯消去）直接判定，比兩者都快；這裡的比較只針對「一般集合」的通用方法。(iv) 時間會隨機器負載波動（見第~\ref{sec:meas} 節開頭的說明）；結論依賴的是數量級（10 倍以上）的差異，而不是個位數百分比。
\end{itemize}

%% =====================================================================
\section{限制與未解之處}

\begin{itemize}[leftmargin=1.5em]
\item 後綴像沒有多項式界；後綴像在 $(\sigma,\lambda)$ 空間中不具結構的集合仍可能很昂貴。量測顯示了哪些情形不會如此。
\item DP 沿用\emph{輸入}決策圖的層序。更好的輸入順序能帶來多大幫助尚未研究。
\item 重排序在 DP 期間只是暫停，沒有被利用。
\item 集合的各層以「每層一個節點清單」保存；若某個集合的決策圖大量跳層，$(u,\ell)$ 配對數 $M\le\lvert f\rvert+\skipc(f)$ 就決定了向下掃描所需的記憶體。
\item 當錯誤集合小到可以列出時，對每個元素直接雜湊 $(\sigma,\lambda)$ 比任何決策圖方法都便宜（表 (c)）。決策圖方法是為了大到無法列出的集合而設計的。
\end{itemize}

\end{document}
